\documentclass{article}
\usepackage [russian] {babel}
\usepackage {mathtext, amsmath, caption, geometry, amssymb, graphicx}
\captionsetup[table]{
	justification=raggedleft,
	singlelinecheck=false,
	labelsep=colon
}
\title{Лабораторная работа №0.0.0. \\ \textbf{Название работы}}
\author{Павел Рудачихин, Б04-507}
\date{00.00.0000}
\begin {document}
\maketitle{}


\section {Аннотация} 
	
	Цель работы + способ. Синтаксис: подтверждение/опровержение гипотезы ... , измерение ... , проверка зависимости .... . 
	

\section {Теоретическое введение}

	Определения. Переход от измеряемой величины к способу измерения. Законы и следствия для рассматриваемых случаев. Без выкладок.
	
	
\section{Описание установки}

	Само собой.
	
	
	
\section{Результаты и обсуждения}

	Графики и пояснения. Все источники погрешности и их значения.
	

\section{Выводы}

	Пересказ обсуждений, может быть даже без чисел. Основной источник погрешности.   


\section{Приложение}

	Таблицы


\end {document}